\documentclass[11pt]{article}

\usepackage[margin=1in]{geometry}
\usepackage{amsmath,amssymb,amsthm}
\usepackage{enumitem}
\usepackage{microtype}
\usepackage{xcolor}
\usepackage[colorlinks=true,linkcolor=blue!50!black,citecolor=blue!50!black,urlcolor=blue!50!black]{hyperref}
\Urlmuskip=0mu plus 1mu
\tolerance=1200
\emergencystretch=3em

\theoremstyle{plain}
\newtheorem{theorem}{Theorem}[section]
\newtheorem{corollary}[theorem]{Corollary}
\newtheorem{lemma}[theorem]{Lemma}
\newtheorem{proposition}[theorem]{Proposition}
\newtheorem{fact}[theorem]{Fact}
\theoremstyle{definition}
\newtheorem{definition}[theorem]{Definition}
\theoremstyle{remark}
\newtheorem{remark}[theorem]{Remark}
\newtheorem{correction}[theorem]{Correction}

\newcommand{\PSL}{\mathrm{PSL}_2(\mathbb{F}_{11})}
\newcommand{\SLg}{\mathrm{SL}_2(\mathbb{F}_{11})}
\newcommand{\edC}{\mathrm{ed}_{\mathbb C}}
\newcommand{\Grr}{\mathrm{Gr}}
\DeclareMathOperator{\Ann}{Ann}
\newcommand{\PP}{\mathbb P}
\newcommand{\dashto}{\dashrightarrow}

\title{Non-Weak-Versality of the Genus-8 Fano Threefold $V_{14}$\\
under $\PSL$: The Centralizer Obstruction}
\author{Transcribed from the E--Klein-Cubic project derivation record}
\date{August 8, 2026}

\begin{document}
\maketitle

\begin{abstract}
We give a self-contained account of the ``centralizer obstruction,'' a
theorem proved in the derivation notes of an ongoing investigation into
the essential dimension of $\PSL$, and of its application to the second
of the two Prokhorov classes of rationally connected $\PSL$-threefolds,
the genus-$8$ Fano threefold $V_{14}$. We state and prove, transcribing
faithfully from the source, a general theorem: if $\sigma$ is an
involution in a finite group $G$ acting faithfully on a smooth projective
variety $Y$ in characteristic $0$, if no faithful linear representation
of $G$ sends $\sigma$ to $\pm\mathrm{id}$ (automatic when $G$ is
centerless), if $N=C_G(\sigma)$, and if (a) no positive-dimensional
irreducible component of $Y^\sigma$ contains a rational curve, and (b)
$Y^N=\emptyset$, then $Y$ is not weakly versal: no $G$-equivariant
rational map from any faithful linear representation of $G$ lands in
$Y$. We then record the machine-verified instantiation of hypotheses
(a) and (b) for $Y=V_{14}$, $G=\PSL$, and $\sigma$ an involution: the
centralizer $N=C_G(\sigma)$ is the dihedral group of order $12$;
$V_{14}^\sigma$ is a single smooth irreducible genus-$1$ sextic curve
together with two reduced points, hence contains no rational curve; and
$V_{14}^{N}$ is empty. Consequently $V_{14}$ is not weakly versal, its
generic twist has no rational point, and $V_{14}$ is not
$\PSL$-unirational --- new-theorem territory, since the standard toolbox
(fixed-point projection, index-two chords, invariant-hyperplane
induction) provably fails on $V_{14}$, and the Cheltsov--Tschinkel--Zhang
equivariant-unirationality program is explicitly scoped to Fano index
$\ge 2$, while $V_{14}$ has index $1$. Combined with Prokhorov's theorem
that the Klein cubic and $V_{14}$ are the only two equivariant
birational classes of rationally connected $\PSL$-threefolds, and with
Duncan--Reichstein's versality theory for invariant cubic hypersurfaces,
this collapses a disjunction: $\edC(\PSL)=3$ if and only if the Klein
cubic itself is $G$-unirational. Whether the Klein cubic is
$G$-unirational remains open. We report exactly what is proved here,
what is cited from the literature, and what is machine-verified, and we
record precisely where each dependency of the argument stands.
\end{abstract}

\section{Introduction}

\subsection{Essential dimension and the versality dictionary}

Let $G$ be a finite group. A \emph{versal $G$-variety} is, informally, a
$G$-variety that receives an equivariant rational map from every
sufficiently generic $G$-variety of the same kind; the precise
definitions are recalled in Section~\ref{sec:prelim}. The
\emph{essential dimension} $\mathrm{ed}(G)$ is the minimal dimension of
a versal $G$-variety carrying a faithful (equivalently, generically
free, for finite $G$) $G$-action~\cite{DR}. Over the base field
$\mathbb C$ we write $\edC(G)$.

Duncan and Reichstein~\cite{DR} organize several competing notions of
versality --- \emph{weakly versal}, \emph{versal}, \emph{very versal}
--- and prove that they correspond exactly to statements about rational
points on twisted forms:
\[
\text{stably birationally linear} \;\Longrightarrow\; \text{very versal}
\;\Longrightarrow\; \text{versal} \;\Longrightarrow\; \text{weakly
versal},
\]
with weak versality equivalent to the existence, for every twisting pair
$(T,K)$, of a $K$-point on the twist ${}^TX$; very versality equivalent
to the existence of a linear representation $G\to GL(V)$ and a
$G$-equivariant \emph{dominant} rational map $V\dashto X$. In the later
terminology of Cheltsov--Tschinkel--Zhang~\cite{CTZ}, very versality of
a projective $G$-variety $X$ is called \emph{$G$-unirationality}:
existence of a dominant equivariant rational map $\PP(V)\dashto X$ for
some $G$-representation $V$. The precise statements are recalled in
Section~\ref{sec:prelim}.

The group $\PSL$ is the automorphism group of the \emph{Klein cubic
threefold}
\[
Y_{\mathrm{Klein}} = \{\, x_0^2x_1+x_1^2x_2+x_2^2x_3+x_3^2x_4+x_4^2x_0=0
\,\}\subset \PP^4,
\]
acting via a linear representation $G\to GL_5(\mathbb C)$~\cite{Adler}.
Since $G$ is simple, the tautological projection $\mathbb C^5\setminus\{0\}
\to \PP^4$ exhibits $\PP^4$ (hence $Y_{\mathrm{Klein}}$, once one uses
that a smooth cubic hypersurface with a rational point is unirational)
as very versal for $G$, giving $\edC(\PSL)\le 4$; and $\PSL$ cannot act
faithfully on a unirational surface, giving $\edC(\PSL)\ge 3$
\cite[Prop.\ 10.8]{DR}. So the headline question of this circle of ideas
is exactly
\[
\edC(\PSL) \in \{3,4\},
\]
and Duncan--Reichstein already frame it as a disjunction of open
conjectures: their Proposition 10.8(b) shows the Cassels--Swinnerton-Dyer
conjecture on $0$-cycles of degree prime to $3$ on cubic hypersurfaces
would force $\edC(\PSL)=3$, while part (c) shows Dolgachev's conjecture
$\mathrm{ed}\ge\mathrm{Crdim}$ would force $\edC(\PSL)=4$ (via
Prokhorov's theorem that $\PSL$ admits no faithful action on any
rational complex threefold). Settling the headline would also complete
Beauville's classification of finite simple groups of essential
dimension $3$~\cite{Beauville}. This note bears on neither conjecture
directly; its role is to narrow \emph{which variety} the question is
really about.

\subsection{The two Prokhorov classes, and where \texorpdfstring{$V_{14}$}{V14} enters}

Prokhorov's classification of finite simple subgroups of the rank-$3$
Cremona group contains the following fact, central to everything below.

\begin{fact}[Prokhorov's two-class theorem {\cite[Thm.\ 1.5]{Prokhorov}}]
\label{fact:prokhorov}
There are exactly two equivariant birational equivalence classes of
rationally connected complex $\PSL$-threefolds, represented by the
Klein cubic threefold $Y_{\mathrm{Klein}}$ and by a smooth Fano
threefold $V_{14}$ of Picard rank $1$, index $1$, and genus $8$ (degree
$14$).
\end{fact}

Consequently, if $\edC(\PSL)=3$, the minimal-dimension versal witness is
(after the standard reduction: a versal variety of minimal dimension is
very versal, hence unirational, hence rationally connected) equivariantly
birational to one of these two classes, and \emph{that} class must
itself be $G$-unirational. The question ``is $\edC(\PSL)=3$?'' therefore
splits into two sub-questions, one about each class. The Klein cubic's
own $G$-unirationality is untouched by this note and remains open. The
present note settles the $V_{14}$ sub-question: $V_{14}$ is \emph{not}
$G$-unirational, and more --- not even weakly versal. This removes
$V_{14}$ from the disjunction entirely, leaving the Klein cubic as the
sole possible witness to $\edC(\PSL)=3$ (Corollary~\ref{cor:ix2} below,
the collapse of Duncan--Reichstein's Remark 10.10).

It is worth recording immediately why this needed a new argument.
Tschinkel--Zhang~\cite{TZ} show that the classical (non-equivariant)
Pfaffian--Grassmannian correspondence between $Y_{\mathrm{Klein}}$ and
$V_{14}$ upgrades, after stabilizing by $\PP^2\times\PP(V)$ for the
irreducible $6$-dimensional representation $V$ of the double cover
$\widetilde G=\SLg$, to a $G$-equivariant birational equivalence; but
the two actions are separately birationally rigid \cite[Thm.\
4.3]{BCDP}, so they are \emph{not} equivariantly birational to each
other, and plain $G$-unirationality of one does not automatically
transfer to the other. The three standard tools that close most
comparable cases in the literature --- projection from a $G$-fixed
point, chords through the fixed pair of an index-$2$ subgroup, and
induction from a $G$-invariant hyperplane section whose induced action
is $G$-unirational \cite[Prop.\ 3.5]{CTZ} --- all fail structurally on
$V_{14}$: the first because the relevant ambient representation is
irreducible (no fixed points), the second because $G$ is simple (no
index-$2$ subgroups), and the third both for the full group (no
invariant hyperplane) and, more strongly, for \emph{every} subgroup
whatsoever, because $V_{14}$ has Fano index $1$ and its hyperplane
sections are K3 surfaces, never unirational. Direct computation with
several subgroups (the $A_5$-fixed locus, the $D_{12}$-fixed pencil)
confirms that the usual fixed-point tools cannot close any $V_{14}$
case either. And Cheltsov--Tschinkel--Zhang's own equivariant
unirationality program~\cite{CTZ} is explicitly scoped to Fano threefolds
of index $\ge 2$; $V_{14}$, at index $1$, falls outside it entirely.
Nothing in the literature closes the $V_{14}$ case. What follows is a
different kind of obstruction --- a \emph{necessary} condition for weak
versality that requires no dominance, no covariant, and no degree data,
checked on a single involution.

\subsection{What is proved here}

Section~\ref{sec:prelim} recalls the precise definitions (versality,
$G$-unirationality, rationally chain connected varieties, fixed loci)
and the three results from the project's equivariant fixed-locus
calculus that the main proof invokes. Section~\ref{sec:main} states and
proves the centralizer obstruction (Theorem~\ref{thm:main} below, called
Corollary IX.1 in the source project notes, where it appears as a
variant of an earlier ``central obstruction'' with the centralizer of an
involution replacing the center of the group). Section~\ref{sec:v14}
records the three inputs needed to instantiate the theorem on $V_{14}$,
all machine-verified. Section~\ref{sec:machine} summarizes the
verification itself: two independent computer-algebra engines, three
primes, and an exact characteristic-$0$ computation. Section~\ref{sec:cor}
draws the conclusion that $V_{14}$ is not weakly versal, with a remark
on what this does and does not say about $V_{14}$ as a bare variety.
Section~\ref{sec:ix2} proves that this collapses the essential-dimension
question onto the Klein cubic alone. Section~\ref{sec:status} is a
frank accounting of what this note actually establishes and what it
depends on. We emphasize at the outset: \textbf{the Klein cubic's own
$G$-unirationality is open}; nothing here proves or suggests a value for
$\edC(\PSL)$, and the note makes no claim about $F_{55}$ or any other
subgroup beyond what is stated.

\section{Preliminaries}
\label{sec:prelim}

Throughout, the base field is $\mathbb C$ (or any algebraically closed
field of characteristic $0$), and $G$ is a finite group.

\begin{definition}[$G$-variety, equivariant rational map]
A \emph{$G$-variety} is a smooth projective integral variety $X$ with a
regular faithful $G$-action. For $H\le G$, $X^H=\{x\in X: hx=x\ \forall
h\in H\}$; in characteristic $0$, with $X$ smooth, $X^H$ is smooth (by
linearizing the $H$-action at a fixed point), so its irreducible
components coincide with its connected components and are pairwise
disjoint for fixed $H$. A rational map $f:X\dashto Y$ of $G$-varieties
is \emph{$G$-equivariant} if it commutes with the two $G$-actions
wherever it is defined; it need not be dominant unless stated.
\end{definition}

\begin{definition}[Versality, following Duncan--Reichstein~\cite{DR}]
Let $G$ be a finite (or, more generally, linear algebraic) group over a
field $k$. A $G$-action on an irreducible $k$-variety $X$ is
\begin{itemize}
\item \emph{weakly versal} if for every field $K/k$ with $K$ infinite
and every $G$-torsor $T\to\operatorname{Spec}(K)$, there is a
$G$-equivariant $k$-morphism $T\to X$;
\item \emph{versal} if every $G$-invariant dense open subvariety of $X$
is weakly versal;
\item \emph{very versal} if there exist a linear representation $G\to
GL(V)$ and a $G$-equivariant dominant rational map $V\dashto X$.
\end{itemize}
Writing $(T,K)$ for a twisting pair and ${}^TX$ for the twist of $X_K$
by $T$, Duncan--Reichstein's Theorem~1.1 identifies these notions with
statements about rational points: $X$ is weakly versal iff ${}^TX(K)
\ne\emptyset$ for every twisting pair $(T,K)$; $X$ is versal iff
$K$-points are dense in ${}^TX$ for every twisting pair; $X$ is very
versal iff ${}^TX$ is $K$-unirational for every twisting pair. An
equivalent, and often more convenient, characterization of weak
versality (\cite[Remark~2.1]{DR}) is: $X$ is weakly versal if and only
if for \emph{every} generically free primitive $G$-variety $Z$ with
$k(Z)^G$ infinite, there is a $G$-equivariant rational map $Z\dashto X$.
Since every faithful linear representation of a \emph{finite} group is
automatically generically free (the union of the proper fixed
subspaces $V^h$, $h\ne 1$, is a proper closed subset, because
faithfulness forces each $V^h\subsetneq V$), a faithful linear $G$-rep
$V$ is itself an instance of such a $Z$; consequently, if $X$ fails to
receive an equivariant rational map from even one faithful linear
representation of $G$, then $X$ is not weakly versal. This is the
precise sense in which the source's informal shorthand ``the generic
twist of $X$ has no rational point'' should be read: by Theorem~1.1(a)
it packages the assertion that \emph{some} twisting pair $(T,K)$
produces a twist ${}^TX$ with ${}^TX(K)=\emptyset$.
\end{definition}

\begin{definition}[$G$-unirationality~\cite{CTZ}]
A $G$-variety $X$ is \emph{$G$-unirational} if there is a dominant
$G$-equivariant rational map $\PP(V)\dashto X$ for some
$G$-representation $V$. Replacing $\PP(V)$ by $V$ shows this notion
coincides with very versality in the sense above.
\end{definition}

\begin{definition}[Rationally chain connected]
A projective variety $Z$ is \emph{rationally chain connected} (RCC) if
two general closed points of $Z$ can be joined by a connected chain of
irreducible rational curves on $Z$. RCC is a birational invariant among
proper varieties, and a projective bundle over an RCC base is RCC.
\end{definition}

We will also need three results from the project's equivariant
fixed-locus calculus~\cite{BComplex}, a working framework built to
track how the fixed-locus structure of a $G$-variety transforms under
equivariant blowups. We state them exactly as needed for the proof in
Section~\ref{sec:main}, with pointers to their labels in the source.

\begin{fact}[Blowup calculus, condensed; {\cite[Thm.\ 2.1]{BComplex}}]
\label{fact:blowup}
Let $Z\subset X$ be a smooth $G$-stable center of a blowup
$\pi:X'=\mathrm{Bl}_ZX\to X$ with exceptional divisor
$E=\PP(N_{Z/X})$. Fix $H\le G$ and an irreducible component $F$ of
$X^H$.
\begin{enumerate}[label=(\roman*)]
\item If $F\not\subset Z$: the strict transform $F'\subset X'$ (the
blowup of $F$ along $F\cap Z$) is again an irreducible component of
$(X')^H$, and is equivariantly birational to $F$.
\item If $F\subseteq Z$: $H$ acts on the normal bundle
$N_{Z/X}|_F$ covering the identity on $F$; write $N_{Z/X}|_F =
\bigoplus_\chi N^\chi$ for the decomposition into $H$-eigen-subbundles
(of rank locally constant along $F$, hence constant since $F$ is
irreducible). For each nontrivial character $\chi$, the closure
$\PP(N^\chi)\subseteq E$ is an irreducible component of $(X')^H$, a
projective subbundle of $E$ over $F$.
\end{enumerate}
\end{fact}

\begin{lemma}[Rational-chain going-down; {\cite[Lem.\ 4.2]{BComplex}}]
\label{lem:goingdown}
Let $q:X\to Y$ be a proper $G$-morphism and $F\subseteq X$ an
irreducible closed RCC subvariety fixed pointwise by some $H\le G$.
Then $q(F)$ is an irreducible closed RCC subvariety of $Y^H$. In
particular, if the irreducible component of $Y^H$ containing $q(F)$
contains no rational curve (e.g.\ is a point, or a curve of geometric
genus $\ge 1$), then $q(F)$ is a single point of it.
\end{lemma}

\begin{proof}
Chains of rational curves push forward under $q$ to chains of rational
curves or points, so a connecting chain between general points of $F$
maps to a connecting chain between general points of $q(F)$; this
shows $q(F)$ is RCC. Pointwise $H$-fixedness of $F$ gives $q(F)\subseteq
Y^H$ by equivariance. If the component of $Y^H$ containing $q(F)$
contains no rational curve, an RCC subvariety of it cannot have
positive dimension (a positive-dimensional RCC variety is covered by
rational curves), so $q(F)$ is a point.
\end{proof}

\begin{lemma}[RCC propagation through the calculus; {\cite[Lem.\
4.3]{BComplex}}] \label{lem:rccprop} In the setting of
Fact~\ref{fact:blowup}, every exceptional stratum $\PP(N^\chi)$ is a
projectivized bundle over $F$ (there written $F_Z$), hence RCC whenever
$F$ is; and strict transforms preserve RCC-ness (being equivariantly
birational to $F$).
\end{lemma}

\begin{correction}[Scope of Lemma~\ref{lem:rccprop}, recorded in the
source as ``Correction I-C,'' 2026-08-05]
\label{corr:ic}
Lemma~\ref{lem:rccprop} is stated above exactly as it should be used:
\emph{a single} exceptional stratum, or strict transform, over \emph{a
single already-RCC} $F$, remains RCC after one blowup. The source
project's notes record that a stronger global consequence once drawn
from it --- that \emph{every} stratum of the fixed-locus decomposition
of \emph{every} $G$-model of $X$ is RCC, as soon as this holds on one
model --- is \textbf{false} for arbitrary towers of equivariant
blowups. Two counterexamples are recorded: one due to A.\ Duncan
(blowing up an isolated fixed point with weights $(-1,-1,-1)$, then a
smooth plane quartic curve inside the resulting exceptional $\PP^2$,
produces a genus-$3$ fixed component on that model); and one found
in-house (blowing up the $G$-orbit of a generic quartic curve inside a
fixed hyperplane of $\PP(W)$ itself produces, over that curve, an
exceptional locus whose fixed part has a section isomorphic to the
quartic, again genus $3$). The per-blowup statement of
Lemma~\ref{lem:rccprop} is unaffected by these examples; what fails is
only the inductive closure of that statement across \emph{arbitrary}
models. The induction does close --- and the stronger global statement
is recovered --- when the blowup centers used are (unions of
components of) strata of the running model, so-called
\emph{stabilizer-stratified towers}. We record this correction because
it is exactly the kind of scope restriction that must be checked
against any argument citing Lemma~\ref{lem:rccprop} across a multi-step
induction; see Remark~\ref{rem:correction-check} below for how the
proof of Theorem~\ref{thm:main} fares against it.
\end{correction}

\begin{fact}[Equivariant resolution / weak factorization; {\cite[Prop.\
3.3]{BComplex}}] \label{fact:resolution} In characteristic $0$, any
$G$-equivariant proper birational morphism (in particular, the
resolution of the graph of a $G$-equivariant rational map $X\dashto Y$
of $G$-varieties) can be dominated by a tower of blowups of smooth
projective $G$-varieties along smooth $G$-invariant centers, resolving
it to a morphism. This rests on Hironaka's resolution of singularities
made canonical (hence automatically equivariant) in characteristic $0$,
together with the equivariant form of the Abramovich--Karu--Matsuki--
Wlodarczyk weak factorization theorem, which the source cites under the
abbreviation ``AKMW.'' Concretely: if $f:X\dashto Y$ is $G$-equivariant
and $\Gamma\subseteq X\times Y$ is the closure of its graph with the
diagonal $G$-action, an equivariant resolution $\widetilde\Gamma\to
\Gamma$ has first projection $\widetilde\pi:\widetilde\Gamma\to X$ a
tower of blowups along smooth $G$-invariant centers, and second
projection $\widetilde f:\widetilde\Gamma\to Y$ an honest (everywhere
defined) $G$-morphism extending $f$.
\end{fact}

Finally, we fix notation for fixed loci used throughout: for $H\le G$
and a $G$-variety $Y$, $Y^H$ denotes the (smooth, by the definition
above) fixed locus; $Y^\sigma:=Y^{\langle\sigma\rangle}$ for an element
$\sigma$; and for a subgroup $N\le G$, $Y^N$ is the locus fixed
pointwise by all of $N$. Since $N\le G$ implies $Y^G\subseteq Y^N$
(anything fixed by all of $G$ is in particular fixed by all of $N$),
$Y^N=\emptyset$ is a strictly stronger statement than $Y^G=\emptyset$.

\section{The centralizer obstruction}
\label{sec:main}

The theorem below is Corollary IX.1 of the source project notes. It is
recorded there under that name because it descends from an earlier
result in the same project, the ``central obstruction'' (Corollary
T3.1): the special case where $\sigma=z$ is taken to be a nontrivial
\emph{central} element of $G$, where $N=C_G(z)=G$ automatically, and
where hypothesis (b) reads $Y^G=\emptyset$. Corollary T3.1 was designed
for groups with nontrivial center; Theorem~\ref{thm:main} below is the
variant anticipated for \emph{local} stabilizers, replacing the center
by the centralizer of a single involution. The trade-off is that the
eigenspace splitting of a representation is only stable under $N=C_G
(\sigma)$ rather than under all of $G$, so the conclusion is
correspondingly narrower: it applies specifically to \emph{linear}
sources (where a single global $N$-stable eigenspace is available
before any resolution is performed), rather than to arbitrary
rationally connected sources carrying a $G$-stable fixed stratum on
some model.

\begin{theorem}[Centralizer obstruction]
\label{thm:main}
Let $G$ be a finite group acting faithfully on a smooth projective
variety $Y$ over an algebraically closed field of characteristic $0$,
and let $\sigma\in G$ be an involution such that no faithful linear
representation $\rho:G\to GL(V)$ under consideration degenerates to
$\rho(\sigma)=\pm\mathrm{id}$. (This is automatic whenever $G$ is
centerless: if $\rho(\sigma)=\mathrm{id}$ then $\sigma\in\ker\rho=\{1\}$
by faithfulness, contradicting that $\sigma$ is an involution; and if
$\rho(\sigma)=-\mathrm{id}$, then since $-\mathrm{id}$ is central in
$GL(V)$, for every $g\in G$ we get $\rho(g\sigma g^{-1})=\rho(g)(-
\mathrm{id})\rho(g)^{-1}=-\mathrm{id}=\rho(\sigma)$, so $g\sigma
g^{-1}=\sigma$ by faithfulness, i.e.\ $\sigma\in Z(G)$; if $Z(G)=1$ this
again forces $\sigma=1$, a contradiction.) Set $N=C_G(\sigma)$. Assume:
\begin{enumerate}[label=(\alph*)]
\item no positive-dimensional irreducible component of $Y^\sigma$
contains a rational curve;
\item $Y^N=\emptyset$.
\end{enumerate}
Then for every faithful linear representation $V$ of $G$, there is no
$G$-equivariant rational map $\PP(V)\dashto Y$ --- dominant or not ---
and none $V\dashto Y$ either (the second follows from the first by
embedding $V\subset\PP(V\oplus\mathrm{triv})$, again linear and
faithful). In particular $Y$ is not weakly versal: some twist of $Y$
has no rational point.
\end{theorem}

\begin{proof}
Fix a faithful linear representation $\rho:G\to GL(V)$. Since $\sigma$
is an involution, $\rho(\sigma)^2=\mathrm{id}$, so $\rho(\sigma)$ is
diagonalizable with eigenvalues in $\{+1,-1\}$; by hypothesis it is
neither $+\mathrm{id}$ nor $-\mathrm{id}$, so both eigenspaces $V_+$
and $V_-$ are nonzero, giving $V=V_+\oplus V_-$. Consider $\PP(V_+)
\subset\PP(V)$: it is nonempty (as $V_+\ne 0$), irreducible and
rational (a linear subspace of projective space), pointwise fixed by
$\sigma$ (by definition of the $+1$-eigenspace, so certainly RCC), and
$N$-\emph{stable}: for $n\in N$ and $v\in V_+$, since $n$ commutes with
$\sigma$ in $G$ we have $\sigma\cdot(n\cdot v)=n\cdot(\sigma\cdot v)=n
\cdot v$, so $n\cdot v\in V_+$ too. This is exactly the point at which
the \emph{linearity} of the representation is used: it supplies a
single subspace $V_+$, stable under $N$, valid at every point of
$\PP(V)$ simultaneously, rather than a subvariety that would have to be
located after the fact.

Suppose, for contradiction, that $\varphi:\PP(V)\dashto Y$ is a
$G$-equivariant rational map (we do not assume it is dominant, or even
non-constant). By Fact~\ref{fact:resolution}, resolve the graph of
$\varphi$ equivariantly: there is a tower of blowups $\pi:\widetilde X
\to\PP(V)$ along smooth $G$-invariant centers, together with an honest
$G$-morphism $\widetilde\varphi:\widetilde X\to Y$ extending $\varphi$.

\smallskip
\noindent\textbf{Claim.} $\widetilde X$ carries an irreducible,
$N$-stable, pointwise-$\sigma$-fixed, RCC closed subvariety
$\widetilde F$.

\smallskip
\noindent\textit{Proof of Claim, by induction up the tower.} Set
$F_0=\PP(V_+)\subseteq\PP(V)$, the base of the tower; by the paragraph
above it is irreducible, RCC (being rational), $N$-stable, and
pointwise $\sigma$-fixed. Suppose $F$ has been constructed with these
four properties at some intermediate stage $X_i$ of the tower, and let
$Z\subset X_i$ be the (smooth, $G$-invariant, hence also
$\sigma$-invariant and $N$-invariant) center of the next blowup $X_{i+1}
=\mathrm{Bl}_ZX_i\to X_i$, with exceptional divisor $E=\PP(N_{Z/X_i})$.

\emph{Case $F\not\subset Z$.} Let $F'\subset X_{i+1}$ be the strict
transform of $F$, i.e.\ the blowup of $F$ along $F\cap Z$ (equivalently,
the closure of the preimage of $F\setminus Z$). By
Fact~\ref{fact:blowup}(i), $F'$ is irreducible and equivariantly
birational to $F$, hence RCC by Lemma~\ref{lem:rccprop}. It is
$N$-stable, since $N$ preserves both $F$ (inductive hypothesis) and $Z$
($G$-invariance), hence preserves $F\cap Z$ and the blowup construction.
It is pointwise $\sigma$-fixed: the open dense subset $F\setminus Z$ of
$F'$ (via the isomorphism $F\setminus Z\xrightarrow{\sim}F'\setminus E$)
is pointwise $\sigma$-fixed by the inductive hypothesis, and being
pointwise fixed by a fixed automorphism is a closed condition, so it
persists on the closure $F'$.

\emph{Case $F\subseteq Z$.} Since $\sigma$ fixes $F$ pointwise and
preserves $Z$, it acts on the total space of the normal bundle
$N_{Z/X_i}|_F$ covering the identity map on $F$; being an involution in
characteristic $0$, this action splits the bundle into eigen-subbundles
$N_{Z/X_i}|_F = N_+\oplus N_-$ (ranks constant along the irreducible
base $F$). As $Z$ has positive codimension in $X_i$, at least one of
$N_+,N_-$ is nonzero; choose $\lambda\in\{+,-\}$ with $N_\lambda\ne 0$
(if both are nonzero, take $\lambda=+$, fixing a convention), and set
\[
F' \;=\; \PP(N_\lambda|_F) \;\subset\; E|_F \;\subset\; X_{i+1}.
\]
This is a projective subbundle of $E$ over $F$ (Fact~\ref{fact:blowup}
(ii), applied to the eigen-character $\lambda$), hence irreducible, and
RCC by Lemma~\ref{lem:rccprop} (a projective bundle over the RCC base
$F$). It is pointwise $\sigma$-fixed: a point of $\PP(N_\lambda|_F)$ is
a line in a fiber of $N_\lambda$, and $\sigma$ scales every vector in
that fiber by the fixed scalar $\lambda=\pm1$, hence fixes the
\emph{line} regardless of the sign. It is $N$-stable: $N$ preserves $F$
and $Z$, hence acts on $N_{Z/X_i}|_F$ covering its action on $F$; since
every $n\in N$ commutes with $\sigma$, $n$ carries the $\lambda$-
eigenspace of $\sigma$ at a point $p\in F$ to the $\lambda$-eigenspace
at $n\cdot p$ (exactly the same commuting argument as for $V_+$ above),
so $n$ preserves the subbundle $N_\lambda$ and hence the subvariety
$F'=\PP(N_\lambda|_F)$.

In both cases $F'$ again has all four properties, completing the
induction; after finitely many steps (the tower is finite) we obtain
$\widetilde F\subseteq\widetilde X$ with all four properties, proving
the Claim.

\smallskip
Now apply $\widetilde\varphi$. Since $\widetilde F$ is pointwise
$\sigma$-fixed and $\widetilde\varphi$ is $G$-equivariant, every $x\in
\widetilde F$ satisfies $\widetilde\varphi(x)=\widetilde\varphi(\sigma
x)=\sigma\widetilde\varphi(x)$, so $\widetilde\varphi(\widetilde F)
\subseteq Y^\sigma$. Since $\widetilde\varphi$ is a morphism of
projective varieties (proper) and $\widetilde F$ is irreducible and
closed, $\widetilde\varphi(\widetilde F)$ is irreducible and closed;
since $\widetilde F$ is RCC, Lemma~\ref{lem:goingdown} shows $\widetilde
\varphi(\widetilde F)$ is an RCC closed subvariety of $Y^\sigma$, lying
(by irreducibility, and smoothness of $Y^\sigma$) inside a single
irreducible component of $Y^\sigma$. By hypothesis (a), if that
component has positive dimension it contains no rational curve, so by
Lemma~\ref{lem:goingdown} $\widetilde\varphi(\widetilde F)$ is forced to
be a single point $y\in Y^\sigma$.

Finally, $\widetilde F$ is $N$-stable and $\widetilde\varphi$ is
$G$-equivariant (in particular $N$-equivariant), so for every $n\in N$:
\[
n\cdot y \;=\; n\cdot\widetilde\varphi(\widetilde F) \;=\;
\widetilde\varphi(n\cdot\widetilde F) \;=\; \widetilde\varphi
(\widetilde F) \;=\; \{y\},
\]
using $n\cdot\widetilde F=\widetilde F$ as sets. Hence $y\in Y^N$,
contradicting hypothesis (b), $Y^N=\emptyset$. This contradiction shows
no such $\varphi$ can exist, for any faithful linear $V$, proving the
theorem.
\end{proof}

\begin{remark}[What is not used]
The proof never assumes $\varphi$ is dominant; a would-be
\emph{constant} equivariant map is excluded too, by the same argument
(its image would be a single $G$-fixed point, and $Y^G\subseteq Y^N=
\emptyset$). No covariant or degree data of any kind enters. The source
remarks that this is the ``all-degree, search-free'' obstruction shape
its broader investigation had been pursuing since an earlier stage of
the project (there labeled ``E34''): a single finite verification, on
one involution, ruling out every faithful linear source of every
dimension and every degree at once.
\end{remark}

\begin{remark}[The proof and Correction~\ref{corr:ic}]
\label{rem:correction-check}
It is worth checking explicitly that the induction above survives
Correction~\ref{corr:ic}, since it invokes Lemma~\ref{lem:rccprop}
inside a multi-step induction along a tower that is an arbitrary
equivariant resolution of a graph (not manifestly a stabilizer-
stratified tower in the sense of the correction). On inspection it
does: at each step the induction tracks a \emph{single} stratum $F$,
and in the case $F\subseteq Z$ it applies Lemma~\ref{lem:rccprop} with
the base of the projective bundle equal to $F$ itself, whose RCC-ness
was already established at the previous stage. This is precisely the
per-blowup statement that Correction~\ref{corr:ic} leaves intact; the
proof never needs the stronger (and, in general, false) global claim
that \emph{every} fixed-locus stratum of \emph{every} model is RCC. We
record this because it is exactly the sort of check a citation of a
scope-corrected lemma demands; the source project notes do not make
this observation explicitly at the point where Theorem~\ref{thm:main}
is stated, so we flag it here as an editorial addition rather than a
transcribed claim.
\end{remark}

\section{The \texorpdfstring{$V_{14}$}{V14} instantiation}
\label{sec:v14}

To apply Theorem~\ref{thm:main} to $Y=V_{14}$ and $G=\PSL$, three facts
must be verified: the identification of $N=C_G(\sigma)$, hypothesis (a),
and hypothesis (b). Since $G$ is simple (in particular centerless), the
degeneracy hypothesis on $\rho(\sigma)$ in Theorem~\ref{thm:main} holds
automatically for \emph{every} faithful linear representation of $G$,
so verifying (a) and (b) once, for $V_{14}$ itself, rules out every
linear source simultaneously.

\subsection{The ambient realization of \texorpdfstring{$V_{14}$}{V14}}

The computations below realize $V_{14}$ concretely as a linear section
of a Grassmannian, following the classical Pfaffian--Grassmannian
correspondence (recalled in Section~\ref{sec:cor}). Let $\widetilde
G=\SLg$ be the double cover of $G$, and let $U$ be its irreducible
$6$-dimensional (``even Weil'') representation. The central element
$-1\in\widetilde G$ acts as $-\mathrm{id}$ on $U$, so $U$ does
\emph{not} descend to a linear representation of $G$ itself (only
$\PP(U)$ carries a genuine, projectively linear, $G$-action; this is
the ``spin'' source discussed and explicitly set aside in
Remark~\ref{rem:spin} below). Its exterior square $\Lambda^2U$ is
$15$-dimensional, and since $-1$ acts on $\Lambda^2U$ as $(-1)^2=+1$,
\emph{this} representation is center-trivial and does descend to a
genuine linear representation of $G=\PSL$. It decomposes as
\[
\Lambda^2U \;\cong\; W_5 \oplus M,
\]
with $W_5$ a $5$-dimensional and $M$ a $10$-dimensional (denoted
``$10'$'' in the source, to distinguish it from the other $10$-
dimensional irreducible of $G$) summand. The variety
\[
V_{14} \;:=\; \Grr(2,U)\cap\PP(M) \;\subset\; \PP(M)=\PP^9
\]
is then the $G$-invariant linear section of the Grassmannian
$\Grr(2,6)\subset\PP(\Lambda^2U)=\PP^{14}$ cut by the linear subspace
$\PP(M)$; its expected dimension is $8-5=3$ (the Grassmannian has
dimension $8$, and $\PP(M)$ has codimension $5$ in $\PP^{14}$), and its
degree, once smoothness holds, is $14$ by transversality of a linear
section of the (arithmetically Cohen--Macaulay) Grassmannian.

\subsection{Input 1: \texorpdfstring{$N=C_G(\sigma)=D_{12}$}{N = centralizer of sigma = D12}}

The class of involutions in $G=\PSL$ is a single conjugacy class of
size $55$. Fixing a representative $\sigma$, the centralizer $N=C_G
(\sigma)$ has order $12$, and its structure --- the dihedral group of
order $12$ (the symmetry group of a hexagon; equivalently, $7$ of its
$11$ nonidentity elements are involutions: $6$ ``reflections'' and one
further involution in the order-$6$ cyclic subgroup) --- is a fact about
the abstract group $G$ alone, independent of any variety it acts on.

\subsection{Input 2: hypothesis (a)}

Since $V_{14}\subset\PP(M)$ is $\sigma$-stable (as $\sigma\in G$ acts
linearly on $M$, preserving $V_{14}$ setwise), and $M=M_+\oplus M_-$
splits into $\sigma$-eigenspaces of dimensions $6$ and $4$ respectively,
\[
V_{14}^\sigma \;=\; \bigl(V_{14}\cap\PP(M_+)\bigr) \;\sqcup\;
\bigl(V_{14}\cap\PP(M_-)\bigr).
\]
The sealed computation (Section~\ref{sec:machine}) finds:
\begin{itemize}
\item $V_{14}\cap\PP(M_+)$ is a curve of degree $6$ with Hilbert
polynomial $6i$ (so arithmetic genus $p_a=1$), whose saturated ideal is
radical with a single associated prime (i.e.\ the curve is reduced and
irreducible) and whose Jacobian minors saturate to the irrelevant
ideal, i.e.\ the curve is \emph{smooth}. A smooth irreducible curve of
arithmetic genus $1$ has geometric genus $1$ and is therefore not
rational.
\item $V_{14}\cap\PP(M_-)$ is a reduced, $0$-dimensional scheme of
degree $2$ --- two points.
\end{itemize}
So $V_{14}^\sigma$ consists of one smooth irreducible genus-$1$ sextic
curve together with two reduced points, and \textbf{contains no
rational curve}: hypothesis (a) holds.

\begin{remark}[Why smoothness is load-bearing]
The Hilbert polynomial $6i$ only certifies the \emph{arithmetic} genus
$p_a=1$ of the sextic; it says nothing about whether the curve is
smooth. An arithmetic-genus-$1$ curve with a single node has geometric
genus $0$ --- its normalization is $\PP^1$ --- and is therefore
\emph{rational}. Had the sextic in $V_{14}\cap\PP(M_+)$ turned out to
be such a nodal curve, hypothesis (a) would fail outright (a rational
curve would sit inside a positive-dimensional component of $V_{14}^
\sigma$), and the entire argument of this section would collapse. This
is exactly the trap the source project notes flag explicitly when
setting up the verification: dimension, degree, and Hilbert polynomial
alone do not suffice; the smoothness of the sextic must be checked
directly (via the Jacobian criterion), and it was.
\end{remark}

\subsection{Input 3: hypothesis (b)}

The subgroup $N=D_{12}$ acts on $M$ ($10$-dimensional); a point of
$\PP(M)$ is fixed pointwise by all of $N$ exactly when it spans a line
on which $N$ acts by some (possibly nontrivial) character, so $V_{14}^N$
decomposes across the isotypic pieces of $M|_N$ for the four
$1$-dimensional characters of $D_{12}$ (the two-dimensional
irreducibles of the nonabelian group $D_{12}$ contain no $N$-stable
lines, hence contribute nothing). These four pieces have dimensions
$2,1,1,0$ (the dimension-$2$ piece is the space of $N$-invariants
proper, i.e.\ the trivial-character piece, matching the pencil found in
earlier, exploratory computation on the same group action). The sealed
computation finds $V_{14}\cap\PP(\text{piece})=\emptyset$ in \emph{all
four} pieces (the fourth, being already $0$-dimensional, is vacuously
empty). Equivalently: restricted to the pencil spanned by the trivial-
character piece, the relevant Plücker data has rank $6$ at all but
three points of the pencil, and rank $4$ (never rank $2$, which is what
membership in $\Grr(2,6)$ would additionally require there) at those
three special points. Hence
\[
V_{14}^{D_{12}} \;=\; \emptyset,
\]
and hypothesis (b) holds.

\section{Machine verification}
\label{sec:machine}

The three inputs above were sealed by a director-run verification
(project packet \url{goal_runs_after_c53d89a/FIX_IX_SEAL}, 2026-08-06),
using two independent computational engines --- exact
rational/modular linear algebra in Python, and Macaulay2 for the ideal-
theoretic computations (saturation, primary decomposition markers,
Jacobian criteria, Hilbert polynomials) --- run at two primes, $397$
and $199$, with an independently-written verifier replaying the entire
computation end-to-end at a third, fresh prime, $353$, and finally an
exact characteristic-$0$ computation over the cyclotomic field $K=
\mathbb Q(\zeta_{11})$ (entries of the explicit matrices lie in
$\mathbb Z[\zeta_{11},1/11]$).

\subsection{The group model}

The representation $U$ is constructed explicitly (an ``even-function''
model on $\mathbb F_{11}$ with a normalized Fourier/Weil generator) via
two $6\times6$ matrix generators, called $T_6,S_6$ in the packet. At
every prime and in the exact field $K$, the following were checked and
passed:
\begin{itemize}
\item $\langle T_6,S_6\rangle$ closes to a group of order exactly
$1320=2\cdot 660$, i.e.\ (the double cover) $\SLg$;
\item the Weil normalization $S_6^2=-I$ and (writing $\mathrm{gauss}$
for the quadratic Gauss sum entering the construction) $\mathrm{gauss}
^2=-11$;
\item the \emph{projective} order profile of the $1320$ elements of the
model --- i.e.\ the multiset of orders of their images in
$G=\PSL$ --- is
\[
(\text{order }1: 2),\ (2: 110),\ (3: 220),\ (5: 528),\ (6: 220),\ (11:
240),
\]
summing to $1320$ as required; this matches $G$'s known element-order
spectrum $\{1,2,3,5,6,11\}$.
\end{itemize}
The isotypic projector onto the $10$-dimensional summand $M$ of
$\Lambda^2U$ has rank exactly $10$ and is idempotent on its own column
space; $\Ann(M)\subset\Lambda^4U$ (under the natural pairing $\Lambda^2U
\times\Lambda^4U\to\Lambda^6U\cong k$) has dimension exactly $5$. The
identification $\Lambda^2U\cong W_5\oplus M$, with $M$ the ``$10'$''
summand rather than the other $10$-dimensional irreducible of $G$, was
confirmed by \emph{two independent methods}: character (isotypic-
projector) averaging, and, in the fresh-prime verifier, an independent
trace-sum method that compares the sum of $\mathrm{tr}(\Lambda^2g)$
over each projective-order class of the $1320$-element model against
the banked character table of $\PSL$ for all four candidate
decompositions ($W_5$ or its Galois twin, paired with either
$10$-dimensional irreducible); only the pairing with the ``$10'$''
candidate matches at every class.

\subsection{The fixed loci and the ambient variety}

At each prime and in $K$: the $\sigma$-eigenspace splitting of $M$ has
dimensions $(6,4)$; the sextic curve $V_{14}\cap\PP(M_+)$ has Hilbert
polynomial $6i$, is radical with one associated prime, and is smooth
(Jacobian minors saturate to the irrelevant ideal); $V_{14}\cap\PP(M_-)$
is a reduced degree-$2$ zero-dimensional scheme. (As a small but
genuine cross-check of independence between primes: the two points of
$V_{14}\cap\PP(M_-)$ are individually rational over $\mathbb F_{397}$
and $\mathbb F_{353}$, but only Galois-conjugate over $\mathbb F_{199}$
--- the scheme is degree-$2$ reduced at every prime regardless, which is
all hypothesis (a) requires.) The centralizer $N=C_G(\sigma)$ has order
$12$ with exactly $7$ projective involutions, matching $D_{12}$; the
four linear-character pieces of $M|_N$ have dimensions $(2,1,1,0)$, and
$V_{14}$ meets none of them.

The ambient variety itself was independently confirmed smooth of pure
dimension $3$: a dual system $J$, consisting of the Pfaffian sextic
(``$\mathrm{Pf}_6$'') together with an adjoint-bivector condition on
$\Ann(M)$, is empty at every prime and over $K$. Since any singular or
excess-dimension point of $\Grr(2,U)\cap\PP(M)$ would produce a point of
$J$, emptiness of $J$ proves $V_{14}$ is smooth of pure dimension $3$;
degree $14$ then follows by transversality, and irreducibility from the
Enriques--Severi--Zariski principle for linear sections of the
(projectively normal) Grassmannian. A direct Gr\"obner-basis computation
of $\Grr(2,U)\cap\PP(M)$ cross-checks dimension $3$ and degree $14$
directly, at all three primes. Finally, the cubic $\mathrm{Pf}_6$ on
$\PP(\Ann M)\cong\PP^4$ is $G$-invariant and nonzero at every prime; by
the (separately sealed) uniqueness of the $G$-invariant cubic on the
$5$-dimensional representation $W_5$, the hypersurface $\{\mathrm{Pf}_6
=0\}$ \emph{is} the Klein cubic threefold. This is the Pfaffian-partner
identification underlying the classical correspondence invoked in
Section~\ref{sec:cor}.

\subsection{Status of the seal}

The packet's own exit marker is \texttt{FIX-IX-SEAL-PASS}. All checks
above appear, with explicit \texttt{PASS} markers, in the packet's log
files (\url{results/checks_397.log}, \url{checks_199.log},
\url{checks_353.log}, \url{checks_K.log}, and the per-prime
\url{m2_sigma_*.out}, \url{m2_d12_*.out}, \url{m2_smooth_*.out},
\url{m2_ambient_*.out} Macaulay2 transcripts), and the independent
end-to-end verifier at the fresh prime $353$ reports
\texttt{ALLGREEN}.

\section{Corollary: \texorpdfstring{$V_{14}$}{V14} is not weakly versal}
\label{sec:cor}

\begin{corollary}
\label{cor:v14}
The action of $G=\PSL$ on $V_{14}$ is not weakly versal: no
$G$-equivariant rational map from any faithful linear representation of
$G$ lands in $V_{14}$, and (by the Duncan--Reichstein dictionary,
Theorem~1.1(a) of~\cite{DR}) some twist of $V_{14}$ has no rational
point. In particular $V_{14}$ is not $G$-unirational.
\end{corollary}

\begin{proof}
$G=\PSL$ is simple, hence centerless, so the degeneracy hypothesis of
Theorem~\ref{thm:main} holds automatically for every faithful linear
representation of $G$ (Section~\ref{sec:v14}). Section~\ref{sec:v14}
verifies $N=C_G(\sigma)=D_{12}$ and hypotheses (a) and (b) of
Theorem~\ref{thm:main} for $Y=V_{14}$, at two primes, a fresh verifying
prime, and in exact characteristic $0$ (Section~\ref{sec:machine}).
Theorem~\ref{thm:main} applies directly.
\end{proof}

This is new-theorem territory: as recalled in the introduction,
Cheltsov--Tschinkel--Zhang's equivariant unirationality program is
scoped to Fano threefolds of index $\ge 2$, and $V_{14}$ has index $1$,
so no result in that literature addresses the $V_{14}$ action at all;
and the standard fixed-point toolbox fails on $V_{14}$ structurally, as
recalled in the introduction.

\begin{remark}[The obstruction is purely equivariant]
\label{rem:bare}
It is important to be precise about what Corollary~\ref{cor:v14} does,
and does not, say about $V_{14}$ as a bare variety, forgetting the
$G$-action. By the classical Pfaffian--Grassmannian correspondence
recalled in Section~\ref{sec:machine} (of which the $G$-equivariant
identification $\{\mathrm{Pf}_6=0\}=Y_{\mathrm{Klein}}$ is the
$G$-invariant instance), $V_{14}$ and $Y_{\mathrm{Klein}}$ are
\emph{birational to each other as bare varieties} --- Tschinkel--Zhang
state this explicitly as ``well-known classically''~\cite{TZ}, and it
is the non-equivariant shadow of their Theorem~1.1. Since birational
invariants transfer across this equivalence: $V_{14}$ \emph{is
unirational} as a bare variety (Kollár's theorem \cite{Kollar}, that a
smooth cubic hypersurface of dimension $\ge 2$ with a rational point is
unirational (recalled as \cite[Thm.\ 10.3]{DR}), applied to $Y_{\mathrm
{Klein}}$ and transported across the birational equivalence), in
particular rationally connected
--- consistent with Fact~\ref{fact:prokhorov} quantifying over
``rationally connected $\PSL$-threefolds.'' But $V_{14}$ is \emph{not
rational}: $Y_{\mathrm{Klein}}$, being a smooth cubic threefold, is
never rational, by the Clemens--Griffiths intermediate-Jacobian
obstruction~\cite{CG}, a theorem with no exceptions among smooth cubic
threefolds; rationality is a birational invariant, so $V_{14}$ is not
rational either. This is exactly Duncan--Reichstein's own summary of
the pair: ``they are unirational \dots\ but not rational''
\cite[Remark~10.10]{DR}.

So Corollary~\ref{cor:v14} is emphatically \emph{not} a repackaging of
the Clemens--Griffiths obstruction, and does not assert that $V_{14}$
is rational --- it is not. What Theorem~\ref{thm:main} adds is strictly
finer, and genuinely equivariant: even restricting attention to sources
$\PP(V)$ for faithful linear representations $V$ of $G$ --- sources
that are themselves rational, so that the bare-variety obstruction to
rationality is simply not in play on the source side either --- no
$G$-\emph{equivariant} map to $V_{14}$ exists. Weak versality is a
strictly $G$-equivariant refinement of unirationality: a variety can be
perfectly unirational as a bare variety (as $V_{14}$ is) while its
$G$-action fails to be weakly versal, exactly as Theorem~\ref{thm:main}
shows here. Nothing in this note bears on the (settled, negative)
question of whether $V_{14}$ is rational as a bare variety; the content
is entirely about the nonexistence of \emph{equivariant} maps.
\end{remark}

\begin{remark}[The spin flank is not addressed here]
\label{rem:spin}
Theorem~\ref{thm:main} quantifies only over \emph{linear} representations
of $G$ itself. The $6$-dimensional $U$ used to build the ambient space
in Section~\ref{sec:v14} is not such a representation: $-1\in\widetilde
G=\SLg$ acts on $U$ as $-\mathrm{id}$, so only $\PP(U)$ carries a
(projectively linear) $G$-action, and $U$ itself is not a source
Theorem~\ref{thm:main} rules out. The source project notes examine this
``spin flank'' separately and find that, at the level of a single
involution $\sigma$, it is \emph{not} obstructed by the same argument:
the lift $\widetilde\sigma\in\widetilde G$ has order $4$, and the
$D_{12}$-stable pair of isolated points in $V_{14}^\sigma$ found in
Section~\ref{sec:v14} matches exactly the shape the spin argument would
require rather than forbid. This does not affect
Corollary~\ref{cor:v14} or Corollary~\ref{cor:ix2} below, both of which
quantify only over honest linear representations of $G$; but it does
mean that this note says nothing about whether $\PP(U)$, or other
projectively-linear ``spin'' sources, admit equivariant maps to
$V_{14}$. That question is left open by the source project notes as
well.
\end{remark}

\section{Corollary IX.2: the disjunction collapses}
\label{sec:ix2}

\begin{corollary}[{Corollary IX.2 of the source}]
\label{cor:ix2}
$\edC(\PSL)=3$ if and only if the Klein cubic threefold $Y_{\mathrm
{Klein}}$ is $G$-unirational.
\end{corollary}

This corollary's inputs are, explicitly: Prokhorov's two-class theorem
(Fact~\ref{fact:prokhorov}); Duncan--Reichstein's Theorem~10.5, that for
a finite group $G$ acting on a smooth $G$-invariant cubic hypersurface
$X\subset\PP(V)$ with $\dim V\ge 4$, the three properties very versal,
versal, and weakly versal are \emph{equivalent} (unconditionally, with
no conjectural input needed for just this equivalence); and
Corollary~\ref{cor:v14} of this note.

\begin{proof}
($\Rightarrow$) Suppose $\edC(\PSL)=3$. Then there is a $3$-dimensional
versal $G$-variety; the standard witness realizing the minimal
dimension (a suitable compression of a generically free linear
representation) is unirational, hence rationally connected
(\cite[Remark~2.6]{DR}). By Fact~\ref{fact:prokhorov}, a rationally
connected $\PSL$-threefold is $G$-birational to either $Y_{\mathrm
{Klein}}$ or $V_{14}$; since versality is a $G$-birational invariant,
one of these two threefolds is itself versal. It cannot be $V_{14}$: a
versal variety is in particular weakly versal, and Corollary~\ref{cor:v14}
shows $V_{14}$ is not weakly versal. So $Y_{\mathrm{Klein}}$ is versal,
hence (Duncan--Reichstein Theorem~10.5, since $Y_{\mathrm{Klein}}
\subset\PP^4=\PP(W_5)$ with $\dim W_5=5\ge 4$ is a smooth $G$-invariant
cubic hypersurface) very versal, i.e.\ $G$-unirational.

($\Leftarrow$) The source records the converse direction as established
separately in the project notes (there labeled ``E37''); we do not
reproduce it here, in keeping with this note's brief to transcribe
Note~IX faithfully rather than to re-derive. We note only, for
orientation, that Duncan--Reichstein's own Proposition~10.8 already
records the unconditional bound $3\le\edC(\PSL)\le 4$ (Section~1.1
above), which is visibly compatible with --- though does not by itself
prove --- this direction.
\end{proof}

As the source puts it, the essential-dimension question becomes
\emph{single-target}: Dolgachev's conjectural route to $\edC(\PSL)=4$
is exactly the assertion ``Klein cubic negative,'' and the
Cassels--Swinnerton-Dyer route to $\edC(\PSL)=3$ is exactly ``Klein
cubic positive''; $V_{14}$ can no longer independently supply an
$\edC=3$ witness. \textbf{We emphasize again: whether $Y_{\mathrm
{Klein}}$ is $G$-unirational is open.} This note proves nothing about
it, in either direction.

\section{Status and dependencies}
\label{sec:status}

We record here, as plainly as possible, what this note establishes and
what it rests on, distinguishing throughout between what is proved in
this note (by transcription from the source), what is cited from
published literature, and what is machine-verified.

\begin{itemize}
\item \textbf{Proved here (transcribed).} Theorem~\ref{thm:main} (the
centralizer obstruction) and its proof; Corollary~\ref{cor:v14} and
Corollary~\ref{cor:ix2} and their proofs (with the converse direction of
Corollary~\ref{cor:ix2} explicitly not reproduced, see above). These are
faithful transcriptions of Corollary IX.1 and Corollary IX.2 of the
source project note \texttt{theory/FIX\_IX\_v14.md}, itself marked
``DRAFT-FOR-DERIVATION'' at the top of that document.

\item \textbf{Depends on a project-internal lemma layer, not published
literature.} The proof of Theorem~\ref{thm:main} rests on
Fact~\ref{fact:resolution} (equivariant resolution/weak
factorization), Lemma~\ref{lem:goingdown} (rational-chain going-down),
and Lemma~\ref{lem:rccprop} (RCC propagation, together with its scope
correction, Correction~\ref{corr:ic}), all from the same project's
\texttt{theory/FIX\_I\_bcomplex.md}. That note is itself explicitly
marked, at its own top: ``Status: DRAFT-FOR-DERIVATION --- every
statement here is to be treated as claimed-until-derived-or-checked,''
with an internal acceptance-test suite (its Section~6, tests T1--T5)
serving as its own validation gate. This is derivation-grade project
work, not peer-reviewed or published mathematics; Fact~\ref{fact:resolution}
rests in turn on genuinely standard results (Hironaka resolution made
canonical in characteristic $0$, and the Abramovich--Karu--Matsuki--
Wlodarczyk equivariant weak factorization theorem), but the specific
formulations of Lemma~\ref{lem:goingdown} and Lemma~\ref{lem:rccprop},
and their packaging into Theorem~\ref{thm:main}, have not been
externally reviewed. We verified directly (Remark~\ref{rem:correction-check})
that the one place Lemma~\ref{lem:rccprop} is used inside a nontrivial
induction survives its own recorded scope correction; we did not
independently re-derive Fact~\ref{fact:resolution},
Lemma~\ref{lem:goingdown}, or the unmarked parts of
Lemma~\ref{lem:rccprop} from first principles.

\item \textbf{Machine-verified.} The geometric inputs of
Section~\ref{sec:v14} --- $N=C_G(\sigma)=D_{12}$, hypothesis (a), and
hypothesis (b) --- are verified computationally, as detailed in
Section~\ref{sec:machine}: two independent engines, two primes ($397$,
$199$) plus an independently-coded fresh-prime verifier ($353$), and an
exact characteristic-$0$ computation over $K=\mathbb Q(\zeta_{11})$.
This is strong evidence by the standards of computer-algebra
verification (independent engines, independent prime, exact
characteristic $0$, and a genuinely independent second method for the
$\Lambda^2U=W_5\oplus M$ identification), but it is a director-run
verification internal to the project, not an independently reproduced
or peer-reviewed computation.

\item \textbf{Cited, not reproved.} Fact~\ref{fact:prokhorov}
(Prokhorov's two-class theorem), Duncan--Reichstein's Theorem~10.5 and
the versality dictionary of Theorem~1.1, Kollár's unirationality
theorem for cubic hypersurfaces, the Clemens--Griffiths non-rationality
theorem, and the identification of $G=\PSL$ as the automorphism group
of the Klein cubic via a linear representation~\cite{Adler}, are all
taken from the literature as stated, without independent re-proof. The
Pfaffian--Grassmannian correspondence and the equivariant birationality
of Tschinkel--Zhang~\cite{TZ} are likewise cited, used here only for
the (non-equivariant) fact that $V_{14}$ and $Y_{\mathrm{Klein}}$ are
birational as bare varieties (Remark~\ref{rem:bare}).

\item \textbf{Consequence, for context.} Corollary~\ref{cor:ix2} makes
the essential-dimension question ``$\edC(\PSL)=3$ or $4$?'' single-
target on the Klein cubic threefold: it is now exactly the question of
whether $Y_{\mathrm{Klein}}$ is $G$-unirational. \textbf{This question
is open.} This note does not claim, suggest, or make progress toward a
value of $\edC(\PSL)$; it does not establish $\edC(\PSL)=4$, and it says
nothing whatsoever about $F_{55}$ or any other subgroup of $G$ beyond
$N=D_{12}=C_G(\sigma)$ and the group $G$ itself. Any claim that the
Klein cubic case, or the value of $\edC(\PSL)$, is resolved would go
beyond what is proved here and beyond what the source project notes
claim.
\end{itemize}

\begin{thebibliography}{S9}
\small

\bibitem{Adler} A.~Adler, \emph{On the automorphism group of a certain
cubic threefold}, Amer.\ J.\ Math.\ \textbf{100} (1978), no.~6,
1275--1280.

\bibitem{Beauville} A.~Beauville, \emph{On finite simple groups of
essential dimension~3}, arXiv:1101.1372.

\bibitem{BCDP} J.~Blanc, I.~Cheltsov, A.~Duncan, Yu.~Prokhorov,
\emph{Finite quasisimple groups acting on rationally connected
threefolds}, arXiv:1809.09226. Cited here, via Tschinkel--Zhang
\cite{TZ}, for their Theorem~4.3: birational rigidity of the
$G$-actions on $Y_{\mathrm{Klein}}$ and $V_{14}$.

\bibitem{CG} C.~H.~Clemens, P.~A.~Griffiths, \emph{The intermediate
Jacobian of the cubic threefold}, Ann.\ of Math.\ (2) \textbf{95}
(1972), 281--356. Citation details as given in Duncan--Reichstein's own
bibliography \cite{DR}; not independently consulted here.

\bibitem{CTZ} I.~Cheltsov, Yu.~Tschinkel, Zh.~Zhang, \emph{Equivariant
unirationality of Fano threefolds}, arXiv:2502.19598.

\bibitem{DR} A.~Duncan, Z.~Reichstein, \emph{Versality of algebraic
group actions and rational points on twisted varieties},
arXiv:1109.6093.

\bibitem{Kollar} J.~Koll\'ar, \emph{Unirationality of cubic
hypersurfaces}, J.\ Inst.\ Math.\ Jussieu \textbf{1} (2002), no.~3,
467--476. Citation details as given in Duncan--Reichstein's own
bibliography \cite{DR}; not independently consulted here.

\bibitem{Prokhorov} Yu.~Prokhorov, \emph{Simple finite subgroups of the
Cremona group of rank~3}, J.\ Algebraic Geom.\ \textbf{21} (2012),
no.~3, 563--600.

\bibitem{TZ} Yu.~Tschinkel, Zh.~Zhang, \emph{Stable equivariant
birationalities of cubic and degree~14 Fano threefolds},
arXiv:2409.08392.

\bibitem[S1]{v14note} \emph{In-repo project source.}
\texttt{problems/E-klein-cubic/theory/FIX\_IX\_v14.md}, Section~5
(``The centralizer obstruction --- the keystone, and the linear-source
case''), the primary source for
Theorem~\ref{thm:main}--Corollary~\ref{cor:ix2}; Sections~1--4 and~6
for surrounding context.

\bibitem[S2]{BComplex} \emph{In-repo project source.}
\texttt{problems/E-klein-cubic/theory/FIX\_I\_bcomplex.md}, the source
for Fact~\ref{fact:resolution} (Prop.\ 3.3), Lemma~\ref{lem:goingdown}
(Lemma~4.2), Lemma~\ref{lem:rccprop} (Lemma~4.3), and its scope
correction (``Correction I-C'').

\bibitem[S3]{Tgate} \emph{In-repo project source.}
\texttt{problems/E-klein-cubic/theory/FIX\_T\_gate.md}, the source for
Corollary T3.1 (the ``central obstruction''), the antecedent of
Theorem~\ref{thm:main} discussed at the start of
Section~\ref{sec:main}.

\bibitem[S4]{SEAL} \emph{In-repo machine verification packet.}
\url{problems/E-klein-cubic/goal_runs_after_c53d89a/FIX_IX_SEAL/}
(\texttt{REPORT.md}, \texttt{brief.md}, \texttt{results/},
\texttt{scripts/}, \texttt{verifier.py}), underlying
Section~\ref{sec:machine}.

\end{thebibliography}

\end{document}

